% ===================================================================
%  அட்டவணை எண் 5 — வீடும் அதன் கட்டுமானமும்.
%
%  Traduction, faite en 2026, en tamoul d'aujourd'hui, sur l'ido de
%  texto/io. Regles de la colonne : voir 10-attavanai-01.tex. Le
%  tableau entier est un DIALOGUE et aucun alinea n'est numerote par
%  l'auteur : toutes les cles portent « 01 », le rang suivant la
%  replique, comme en ido.
%
%  LE TABLEAU LE PLUS LONG DU LIVRET — soixante-huit blocs, plus de
%  cent vingt renvois — et le troisieme ecrit avec la liste de
%  parigi.py en main des la premiere ligne, apres le telougou et le
%  coreen. C'est ici que l'ordre des renvois coute le plus cher si on
%  le decouvre apres coup.
%
%  ONZE RETOURNEMENTS SUR TRENTE-SEPT PAIRES. La colonne coreenne en
%  avait compte douze au meme tableau ; l'ecart tient a une seule
%  paire — « la kalko esas proxim ilu en sako (71) » —, que le tamoul
%  laisse en place parce que ses deux adjoints se suivent dans
%  l'ordre de l'ido. C'est le taux le plus BAS de la colonne tamoule
%  jusqu'ici, plus bas que le defile du tableau 3 (dix-sept sur
%  vingt-cinq), et la raison est celle du marche telougou : un
%  chantier s'enumere. On nomme les corps de metier l'un apres
%  l'autre, et les outils de chacun a la suite.
%
%  ET renvoji.py A RELEVE TROIS ECARTS AU PREMIER PASSAGE, dont DEUX
%  QUE parigi.py NE POUVAIT PAS LISTER. Au bloc 01-30 j'avais mis les
%  plafonds (26) avant le platre (70) ; au bloc 01-38 la toile (29)
%  avant l'echelle (106) et le balcon (28). Ni « indutas per gipso la
%  plafoni » ni « pozas storo » n'est un rapport modificateur–tete :
%  parigi.py ne les voit pas, et n'a pas a les voir. Le troisieme
%  ecart etait une phrase laissee a moitie refaite, qui doublait le
%  renvoi (14). Trois fautes d'execution, aucune de doctrine — et
%  c'est exactement le partage du travail entre les deux outils : la
%  liste dit ce qu'il faut prevoir, le controle dit ce qu'on a
%  reellement ecrit. Une colonne qui n'aurait que la liste
%  passerait a cote de deux ecarts sur trois.
%
%  kolonoj.py A RELEVE UNE QUATRIEME CHOSE : un \nl reste dans la
%  note du bloc 01-25, recopie du fac-simile par distraction. Une
%  traduction ne suit aucune ligne imprimee et n'a donc aucun \nl —
%  c'est la premiere des trois regles de cette colonne, et le
%  controle de forme la tient meme quand la main l'oublie.
%
%  LES METIERS DU BATIMENT ONT TOUS LEUR NOM EN TAMOUL, ET ILS SONT
%  VIEUX : கொத்தனார் le macon, தச்சர் le charpentier, கொல்லர் le
%  forgeron, தோட்டக்காரர் le jardinier. Les outils aussi : கரணை la
%  truelle, தூக்குக் குண்டு le fil a plomb, இழைப்புளி le rabot, உளி
%  le ciseau, அடைகல் l'enclume, இடுக்கி les tenailles, அரம் la lime,
%  தாழ்ப்பாள் le verrou, சாரம் l'echafaudage, ஓடு la tuile,
%  பாத்திகள் les planches du jardin. Les emprunts sont pour les
%  materiaux et pour eux seuls — சிலேட்டு l'ardoise, துத்தநாகம் le
%  zinc n'en est pas un, பிளாஸ்டர் le platre, பசை le mastic non plus.
%
%  MAIS LE TAMOUL N'A QU'UN MOT LA OU L'IDO EN A DEUX. தச்சர் couvre
%  a la fois le charpentier (115), celui qui fait la charpente, et le
%  menuisier, celui qui fait les portes, les fenetres, les parquets
%  et les volets. Le coreen distinguait 목수 et 소목장이. La colonne
%  ne forge pas un mot : elle qualifie — கூரைத் தச்சர், le charpentier
%  de toiture, et மரத் தச்சர், le charpentier de bois ouvre. Nommer la
%  difference est honnete ; inventer un mot pour elle ne le serait
%  pas.
%
%  ET DEUX OUTILS SE RESSEMBLENT ASSEZ POUR QU'ON LES SEPARE. Le
%  « shovelo » (82) du terrassier et le « spado » (131) du jardinier
%  sont deux choses differentes, et le tamoul les separe sans effort :
%  மண் அள்ளி, ce qui ramasse la terre, contre மண்வெட்டி, ce qui la
%  coupe. C'est la planche qui tranche, pas le dictionnaire.
%
%  LE BALCON (28) EST UN CINQUIEME PIEGE, ET ON LE REFUSE COMME LES
%  QUATRE AUTRES. Le tamoul dirait volontiers திண்ணை ou மொட்டை
%  மாடம் — la banquette-galerie, la terrasse plate —, mais le premier
%  a deja ete refuse au tableau 4 pour le verando et le second est la
%  terrasse (21) de ce tableau-ci. La colonne ecrit பால்கனி, comme
%  elle a ecrit வராந்தா.
%
%  LE FAC-SIMILE REEMPLOIE TROIS DE SES RENVOIS et la source ne bouge
%  pas : (107) est le porteur de mortier au bloc 01-21 et le
%  tapissier au bloc 01-38 ; (33) est le verando au bloc 01-30 et les
%  mansardes au bloc 01-55 ; (1) est la facade ET la salle de bain
%  dans la meme phrase du bloc 01-53 — parigi.py le signale d'ailleurs
%  en listant « (1) ← di ← (1) ». La traduction porte les six, dans
%  l'ordre recu.
% ===================================================================

%%K t05-apar-1 apar
\VUsaut{6.0mm}
\VUtitre{15.0pt}{60}{அட்டவணை எண் 5}
\VUsaut{1.4mm}
\VUtitre{11.4pt}{0}{வீடும் அதன் கட்டுமானமும்.}
\VUsaut{2.2mm}

%%K t05-tit-1 sub
\VUpk{10.2pt}{40}{நல்ல சந்திப்பு.}
\VUsaut{3.0mm}

%%K t05-01-1 p
\textsc{இயோவான்னெஸ்}. --- மாமா, \VUgras{வீட்டை} எப்படிக் கட்டுகிறார்கள்
என்பதை அறிய நான் மிகவும் விரும்புகிறேன்.

%%K t05-01-2 p
\textsc{மாமா} (80). --- அது மிகவும் எளிது, என் நண்பா ; என்னுடன் வா.
பெர்னார்துஸ் ஐயா (79) கட்டச் சொல்லியிருக்கும், நான் குடியிருக்கப் போகும்
வீட்டை உனக்குக் காட்டுகிறேன்.

%%K t05-01-3 p
\textsc{இயோவான்னெஸ்}. --- இந்த வீட்டின் \VUgras{திட்டத்தை}
\textsuperscript{(76)} வரைந்தது யார்?

%%K t05-01-4 p
\textsc{மாமா}. --- \VUgras{கட்டிடக் கலைஞர்} \textsuperscript{(75)}. அவர்
மிகவும் தெளிவான வரைபடம் ஒன்றைக் கொடுத்தார், அது ஒப்புக் கொள்ளப்பட்டது ;
பிறகு \VUgras{ஒப்பந்தக்காரர்} \textsuperscript{(77)} வேலைகளைத் தொடங்கினார்.

%%K t05-01-5 p
\textsc{இயோவான்னெஸ்}. --- ஒப்பந்தக்காரர் யார்?

%%K t05-01-6 p
\textsc{மாமா}. --- அவர் பிளாங்கோ ஐயா, உன் நண்பன் யாகோபுஸின் தந்தை. அவர்
கட்டிடக் கலைஞர் ரூஃபோ ஐயாவுடன் சேர்ந்து வேலையாட்களை நடத்துகிறார்.

%%K t05-01-7 p
இதோ! சரியான நேரத்தில் பிளாங்கோ ஐயா தன் \VUgras{அளவுகோலுடன்}
\textsuperscript{(78)} வருகிறார் ; ரூஃபோ ஐயாவும் தன் திட்டத்துடன் வருகிறார்.

%%K t05-01-8 p
வணக்கம் ஐயாக்களே. நலமா?

%%K t05-01-9 p
\textsc{ரூஃபோ ஐயா}. --- மிக நலம், நன்றி. வணக்கம் இயோவான்னெஸ், இந்தக்
குழந்தை எவ்வளவு வளர்ந்து விட்டது!

%%K t05-01-10 p
\textsc{மாமா}. --- ஆம், உண்மைதான், இவன் ஒரு சிறு மனிதன். இவன் மிகவும்
நிதானமானவன் ; தான் பார்க்கும் எல்லாவற்றைப் பற்றியும் இவனுக்குச் சொல்ல
வேண்டும். இன்று வீட்டை எப்படிக் கட்டுகிறார்கள் என்பதை அறிய
விரும்புகிறான்.

%%K t05-tit-2 sub
\VUpk{10.2pt}{40}{வேலையாட்கள்.}
\VUsaut{3.0mm}

%%K t05-01-11 p
\textsc{பிளாங்கோ ஐயா}. --- வா என்னுடன், சிறு நண்பா ; உடனே உனக்கு
விளக்குகிறேன், ஏனெனில் கற்க விரும்பும் சிறுவர்களை நான் மிகவும்
விரும்புகிறேன்.

%%K t05-01-12 p
கையில் \VUgras{மண் அள்ளியைப்} \textsuperscript{(82)} பிடித்திருக்கும் அந்த
வேலையாளைப் பார்க்கிறாய் : அவர் \VUgras{மண் வேலைக்காரர்}
\textsuperscript{(81)}. நீர்க் குழாய்களைப் போடுவதற்காக \VUgras{அகழியை}
\textsuperscript{(46)} அவர் தோண்டுகிறார். கொத்தனார் நான்கு \VUgras{சுவர்களை}
எழுப்புவதற்காக, வீடு இருக்கும் இடத்தில் தரையையும் அவர் தோண்டினார். அந்தச்
சுவர்களில் தரைக்குள் இருக்கும் பகுதிக்கு \VUgras{அடித்தளம்}
\textsuperscript{(2)} என்று பெயர். அடித்தளங்களுக்கு இடையில் அடங்கிய
இடம் \VUgras{நிலவறையை} \textsuperscript{(3)} உண்டாக்குகிறது. அந்த
நிலவறைக்கு \VUgras{நிலவறைச் சாளரம்} \textsuperscript{(5)} என்னும் சிறிய
சாளரமும், \VUgras{கதவும்} \textsuperscript{(4)}, படிக்கட்டும் உண்டு.

%%K t05-01-13 p
\VUgras{கொத்தனார்} \textsuperscript{(108)} \VUgras{கதவுகளுக்கும்}
\textsuperscript{(8)} \VUgras{சாளரங்களுக்கும்} \textsuperscript{(17)}
திறப்புகளை விட்டுச் சுவர்களைத் தொடர்ந்து கட்டினார்.

%%K t05-01-14 p
\textsc{இயோவான்னெஸ்}. --- ஆனால் அந்தக் கொத்தனாரை நான் பார்க்கவில்லையே!

%%K t05-01-15 p
\textsc{பிளாங்கோ ஐயா}. --- இப்போது அவர் வீட்டில் வேலை செய்து முடித்து
விட்டார். அவர் \VUgras{குதிரை லாயத்திற்கு} \textsuperscript{(54)} அருகில்,
தன் \VUgras{சாரத்தின்} \textsuperscript{(110)} மேல் இருக்கிறார் ; அவர்
\VUgras{துணி துவைக்கும் அறையின்} \textsuperscript{(56)} சுவரைக் கட்டுகிறார்.

%%K t05-01-16 p
அவரிடம் கரணையும், சாந்துத் தட்டும், \VUgras{தூக்குக் குண்டும்}
\textsuperscript{(109)} இருக்கின்றன. ஆனால் இந்தப் பெயர்கள் உனக்கு நினைவில்
இருக்காது.

%%K t05-01-17 p
\textsc{இயோவான்னெஸ்}. --- நிச்சயம் இருக்கும் ; ஏனெனில் உடனே என் மாமாவிடம்
சிறு கரணையும் சாந்துத் தட்டும் கேட்கப் போகிறேன் ; தூக்குக் குண்டை நானே
ஒரு சிறு கயிற்றால் செய்து கொள்வேன்.

%%K t05-tit-3 sub
\VUsaut{3.0mm}
\VUpk{10.2pt}{40}{கட்டுமானப் பொருள்.}
\VUsaut{3.0mm}

%%K t05-01-18 p
\textsc{மாமா}. --- ஆ! மிக நல்லது. ஆனால் சொல் இயோவான்னெஸ், வீடு கட்ட
என்ன வேண்டும்?

%%K t05-01-19 p
\textsc{இயோவான்னெஸ்}. --- \VUgras{செதுக்கிய கற்களும்} \textsuperscript{(59)}
\VUgras{சாந்தும்} \textsuperscript{(65)} வேண்டும்.

%%K t05-01-20 p
\textsc{மாமா}. --- சரிதான். பெரிய தொப்பியுடன் தன் \VUgras{வண்டியின்}
\textsuperscript{(136)} மேல் இருக்கும் அந்த மனிதரைப் பார். அவர்
\VUgras{வண்டிக்காரர்} \textsuperscript{(135)}. நாள்தோறும் \VUgras{கற்பாளங்களை}
\textsuperscript{(60)} இங்கே கொண்டு வருகிறார். முற்றத்தில் தன் வண்டியை
இறக்குகிறார் ; \VUgras{கல் செதுக்குபவர்கள்} \textsuperscript{(83)} அந்தப்
பாளங்களைச் செதுக்கி, தங்கள் \VUgras{கல் அறுவாளால்} \textsuperscript{(85)}
அறுக்கிறார்கள். \VUgras{உதவி ஆள்} \textsuperscript{(146)} அவற்றைக்
கொத்தனாரிடம் கொண்டு போகிறார் ; கொத்தனார் அவற்றை இடத்தில் வைக்கிறார்.

%%K t05-01-21 p
அதே வேளையில் \VUgras{சாந்து கலப்பவர்} \textsuperscript{(87)} சாந்தைத் தயார்
செய்கிறார். அவருக்கு \VUgras{மணலும்} \textsuperscript{(62)},
\VUgras{சுண்ணாம்பும்} \textsuperscript{(63)}, \VUgras{நீரும்}
\textsuperscript{(64)} வேண்டும். சுண்ணாம்பு அவருக்கு அருகில் ஒரு
\VUgras{சாக்கில்} \textsuperscript{(71)} இருக்கிறது ;
\VUgras{கொத்தனாரின் உதவியாள்} \textsuperscript{(111)} ஒரு
\VUgras{தள்ளுவண்டியில்} \textsuperscript{(112)} மணலைக் கொண்டு வருகிறார் ;
\VUgras{பயிற்சி ஆள்} \textsuperscript{(113)} இறைப்பானின் (58) அருகில்
இருக்கிறார். அவர் ஒரு \VUgras{வாளியை} \textsuperscript{(114)}
நிரப்புகிறார். சாந்து விரைவில் தயாராகி விடுகிறது.
\VUgras{சாந்து சுமப்பவர்} \textsuperscript{(107)} அதை எடுத்து, ஏணி வழியாகச்
சாரத்தின் மேல் தூக்கிச் செல்கிறார் ; அங்கே வீட்டின் அந்தப் பக்கத்தை
முடிக்கும் கொத்தனார் வேலை செய்கிறார்.

%%K t05-01-22 p
இப்போது சொல், \VUgras{கூரையில்} \textsuperscript{(31)} வேலை செய்யும்
வேலையாளுக்கு என்ன பெயர்?

%%K t05-01-23 p
\textsc{இயோவான்னெஸ்}. --- அவர் \VUgras{கூரை வேலைக்காரர்}
\textsuperscript{(118)}.

%%K t05-tit-4 sub
\VUsaut{3.0mm}
\VUpk{10.2pt}{40}{வீட்டின் கூரை.}
\VUsaut{3.0mm}

%%K t05-01-24 p
\textsc{பிளாங்கோ ஐயா}. --- ஆம் ; ஆனால் முதலில் \VUgras{கூரைத் தச்சர்}
\textsuperscript{(115)} வருகிறார்.

%%K t05-01-25 p
கூரைத் தச்சர் \VUgras{மரச் சட்டத்தில்} \textsuperscript{(35)} வேலை
செய்கிறார். அவர் \VUgras{உத்திரங்களை} \textsuperscript{(37)}
\VUgras{மர ஆணிகளால்} \textsuperscript{(117)} இணைக்கிறார். அங்கே மேலே,
அகன்ற வெல்வெட் குட்டைக் கால்சட்டையுடன் இருக்கும் அந்த வேலையாட்கள்
கூரைத் தச்சர்கள். ஒருவர் சிறு உத்திரத்தைப் பிடித்திருக்கிறார் ; இன்னொருவர்
தன் \VUgras{கோடரியால்} \textsuperscript{(116)} ஆணியை வெட்டுகிறார்.
\VUgras{புகைபோக்கிக்கு} \textsuperscript{(41)} அருகில் இருப்பவர் கூரை
வேலைக்காரர். அவர் (*) விட்டங்களின் மேல் பட்டைகளையும், பட்டைகளின் மேல்
\VUgras{சிலேட்டுத் தகடுகளையும்} \textsuperscript{(66)} அறைகிறார். சில
வேளைகளில் ஓடுகளைப் போடுகிறார்.

%%K t05-noto-1 noto
\VUnotes{143.2mm}{%
(*) \VUgras{Balk-o} = ஜெ. \textit{Balken}, ஆ. \textit{joist}, பி.
\textit{solive}, இ. \textit{travicello}, ர. \textit{balka}, ஸ். போ.
\textit{viga}.
இதுவரை ஏற்றுக் கொள்ளப்படாத தொழில்நுட்ப வேர்ச்சொல் ; ஆனால் பி.
\textit{solive} என்னும் சொல்லின் பொருளை மொழிபெயர்க்க முன்மொழியத்
தக்கது. அது trabo (பி. \textit{poutre}) அல்ல, trabeto வும் அல்ல. அது
தளத்தையும் உட்கூரையையும் தாங்கும் சதுரமாகச் செதுக்கிய மரக் கழி.%
}

%%K t05-01-26 p
குதிரை லாயத்தின் கூரை \VUgras{ஓடு வேய்ந்தது} \textsuperscript{(55)},
வீட்டின் கூரை \VUgras{சிலேட்டு வேய்ந்தது} \textsuperscript{(32)},
வராந்தாவின் கூரை \VUgras{துத்தநாகத்தால்} \textsuperscript{(24)} ஆனது.

%%K t05-01-27 p
\textsc{இயோவான்னெஸ்}. --- துத்தநாகக் கூரைகளிலும் கூரை வேலைக்காரர்தான்
வேலை செய்கிறாரா?

%%K t05-01-28 p
\textsc{பிளாங்கோ ஐயா}. --- இல்லை ; \VUgras{துத்தநாக வேலைக்காரர்}
\textsuperscript{(121)} அல்லது \VUgras{ஈய வேலைக்காரர்} செய்கிறார் ;
அவர்தான் வீட்டின் கூரையின் மேல் \VUgras{கூரைச் சாளரத்தை}
\textsuperscript{(38)} வைக்கிறார். \VUgras{மழை நீர்க் குழாய்களையும்}
\textsuperscript{(43)} \VUgras{கூரைக் கால்வாய்களையும்} \textsuperscript{(42)}
அவரே வைக்கிறார். துத்தநாகத்தையும் தகரத்தையும் அவர் ஒட்ட வைக்கிறார்.
\VUgras{மின்னல் தடுப்பையும்} \textsuperscript{(40)} \VUgras{காற்றுக்
காட்டியையும்} \textsuperscript{(39)} அவர் பொருத்தினார் ; அவை
\VUgras{முகட்டுச் சுவரின்} \textsuperscript{(36)} மேல் இருக்கின்றன.

%%K t05-01-29 p
\textsc{இயோவான்னெஸ்}. --- அப்படியானால் வீடு முற்றிலும் முடிந்து விட்டதா?

%%K t05-tit-5 sub
\VUsaut{3.0mm}
\VUpk{10.2pt}{40}{கருவிகள்.}
\VUsaut{3.0mm}

%%K t05-01-30 p
\textsc{பிளாங்கோ ஐயா}. --- முழுவதும் இல்லை. \VUgras{பூச்சு வேலைக்காரர்}
\textsuperscript{(99)} \VUgras{செங்கற்களால்} \textsuperscript{(61)} பிரிச்
சுவர்களைக் கட்டுகிறார் ; அவை வீட்டைப் பல அறைகளாகப் பிரிக்கின்றன.
அவர் \VUgras{பூச்சுச் சாந்தால்} \textsuperscript{(70)}
\VUgras{உட்கூரைகளையும்} \textsuperscript{(26)} சுவர்களையும் பூசுகிறார். இப்போது
\VUgras{வராந்தாவின்} \textsuperscript{(33)} உட்கூரையில் வேலை செய்கிறார்.
கொத்தனாரைப் போலவே அவரிடமும் \VUgras{கரணையும்} \textsuperscript{(100)}
\VUgras{சாந்துத் தட்டும்} \textsuperscript{(101)} இருக்கின்றன.

%%K t05-01-31 p
பிறகு \VUgras{மரத் தச்சர்} கதவுகள், சாளரங்கள், \VUgras{மரத் தளங்கள்}
\textsuperscript{(16)}, \VUgras{சாளரக் கதவுகள்} \textsuperscript{(19)}
ஆகியவற்றில் வேலை செய்கிறார்.

%%K t05-01-32 p
இப்போது அவர் முற்றத்தில் தன் \VUgras{வேலை மேசையின்} \textsuperscript{(90)}
மேல் வேலை செய்கிறார். அவரிடம் \VUgras{அறுவாள்} \textsuperscript{(94)}
இருக்கிறது — அது \VUgras{மரத்தை} (69) அறுக்கிறது —, மேலும்
\VUgras{இழைப்புளி} \textsuperscript{(91)}, \VUgras{இறுக்கி}
\textsuperscript{(92)}, \VUgras{உளி} \textsuperscript{(94 bis)},
\VUgras{கவராயம்} \textsuperscript{(95)}, மரச் \VUgras{சுத்தி}
\textsuperscript{(95 bis)} இருக்கின்றன.

%%K t05-01-33 p
\textsc{இயோவான்னெஸ்}. --- ஆ! கொத்து வேலையை விட மர வேலை மிகவும்
மகிழ்ச்சி தருவது. மாமா, எனக்கு வேலை மேசையும் அறுவாளும் இழைப்புளியும்
வாங்கித் தருவாயா? விரைவில் மேதோருக்கு ஒரு குடிசை செய்யப் போகிறேன்.

%%K t05-01-34 p
\textsc{மாமா}. --- ஆம், சிறுவனே.

%%K t05-01-35 p
\textsc{இயோவான்னெஸ்}. --- நான் எவ்வளவு மகிழ்ச்சியாக இருக்கிறேன்! நுழைவுக்
கதவருகில் நிற்பவர் யார்?

%%K t05-tit-6 sub
\VUsaut{3.0mm}
\VUpk{10.2pt}{40}{வண்ணம் பூசுதல்.}
\VUsaut{3.0mm}

%%K t05-01-36 p
\textsc{பிளாங்கோ ஐயா}. --- அவர் \VUgras{வண்ணம் பூசுபவர்}
\textsuperscript{(102)}. \VUgras{வண்ணப்} \textsuperscript{(103)} பானையுடனும்
தன் \VUgras{தூரிகையுடனும்} \textsuperscript{(104)} தன் ஏணியின் மேல்
நிற்கிறார். நுழைவுக் கதவின் மரம் நெடுங்காலம் நிலைப்பதற்காக அவர் அதற்கு
வண்ணம் பூசுகிறார்.

%%K t05-01-37 p
அறைகளில் காகிதங்களை ஒட்டுவதும் அவரே.

%%K t05-01-38 p
\VUgras{திரைச் சீலை வேலைக்காரர்} \textsuperscript{(107)} — அவர் ஒரு
\VUgras{ஏணியின்} \textsuperscript{(106)} மேல், \VUgras{பால்கனியில்}
\textsuperscript{(28)} இருக்கிறார் — \VUgras{திரைச் சீலையை}
\textsuperscript{(29)} பொருத்துகிறார்.

%%K t05-01-39 p
பெரிய சாளரத்திற்கு அருகில் நிற்கும் வேலையாள் \VUgras{கண்ணாடி
வேலைக்காரர்} \textsuperscript{(96)}. அவர் \VUgras{கண்ணாடித் தகடுகளை}
\textsuperscript{(18)} \VUgras{பசையால்} \textsuperscript{(73)}
பொருத்துகிறார். நிலவறையின் கதவருகில் மண்டியிட்டிருக்கும் அந்த இன்னொரு
வேலையாள் \VUgras{பூட்டு வேலைக்காரர்} \textsuperscript{(97)}. அவர்
\VUgras{பூட்டை} \textsuperscript{(6)} பொருத்துகிறார். அவருடைய
\VUgras{அரம்} \textsuperscript{(98)} அவருக்கு அருகில் இருக்கிறது.

%%K t05-01-40 p
\textsc{இயோவான்னெஸ்}. --- தன் \VUgras{சுத்தியால்} \textsuperscript{(84)}
இரும்பின் மேல் அடிக்கும் அந்த வேலையாளும் பூட்டு வேலைக்காரரா?

%%K t05-01-41 p
\textsc{பிளாங்கோ ஐயா}. --- இன்னும் சரியாகச் சொன்னால், அவர் கொல்லர் (125).
அவர் \VUgras{இரும்புப் பொருள்களை} \textsuperscript{(67)} அடிக்கிறார் ; அவை
\VUgras{மின் வேலைக்காரருக்கு} \textsuperscript{(124)} உரியவை. அந்த மின்
வேலைக்காரர் \VUgras{வண்டிக் கொட்டகையின்} \textsuperscript{(51)} கூரையின்
மேல் இருக்கிறார். கொல்லரிடம் \VUgras{உலையும்} \textsuperscript{(126)},
\VUgras{அடைகல்லும்} \textsuperscript{(127)}, \VUgras{இடுக்கியும்}
\textsuperscript{(128)} இருக்கின்றன.

%%K t05-01-42 p
மின் வேலைக்காரர் தொலைபேசியின் \VUgras{செம்புக்} \textsuperscript{(44)}
\VUgras{கம்பியைப்} பொருத்துகிறார்.

%%K t05-tit-7 sub
\VUpk{10.2pt}{40}{தோட்ட வேலை.}
\VUsaut{3.0mm}

%%K t05-01-43 p
\textsc{மாமா}. --- \VUgras{முற்றத்தின்} \textsuperscript{(45)} உள்பக்கத்தில்,
\VUgras{இரும்புக் கிராதிக்கும்} \textsuperscript{(47)} \VUgras{பெருங்கதவுக்கும்}
\textsuperscript{(48)} அருகில் \VUgras{தோட்டக்காரரைப்} \textsuperscript{(129)}
பார்க்கிறாய். அவர் \VUgras{சிறு தோட்டத்தைத்} \textsuperscript{(49)} தயார்
செய்கிறார். தன் \VUgras{மண்வெட்டியால்} \textsuperscript{(131)} மண்ணைக்
கொத்தினார், விதைகளை விதைக்கிறார், \VUgras{இரும்புச் சீப்பால்}
\textsuperscript{(130)} சமன் செய்கிறார். பிறகு தன் \VUgras{நீர்த்
தெளிப்பானால்} \textsuperscript{(132)} \VUgras{பாத்திகளுக்கு}
\textsuperscript{(150)} நீர் விடுவார், செடிகள் வளரும்.

%%K t05-01-44 p
குதிரையைப் பராமரித்து உணவு கொடுத்து முடித்த \VUgras{குதிரைக் காரர்}
\textsuperscript{(133)} \VUgras{வண்டியைக்} \textsuperscript{(52)}
கடற்பஞ்சாலும், \VUgras{கைப் பம்பாலும்} \textsuperscript{(134)}, ஒரு வாளி
நீராலும் கழுவுகிறார். பிறகு அதை வண்டிக் கொட்டகைக்குள் மீண்டும்
கொண்டு போய், தடித்த \VUgras{தாழ்ப்பாளால்} \textsuperscript{(53)} கதவை
அடைப்பார்.

%%K t05-01-45 p
இதோ நீ மகிழ்ச்சியாக இருக்கிறாய் என்று நினைக்கிறேன் ; போதுமான
விளக்கங்கள் உனக்குக் கிடைத்தன.

%%K t05-01-46 p
\textsc{இயோவான்னெஸ்}. --- ஆம் மாமா ; ஆனால் வீட்டின் உள்ளேயும் நான்
விருப்பத்துடன் பார்ப்பேன்.

%%K t05-01-47 p
\textsc{மாமா}. --- அது நாளை நடக்கும். ரூஃபோ ஐயா உனக்குத் திட்டத்தை
விளக்குவார், ஒவ்வோர் அறையின் பெயரையும் சொல்வார்.

%%K t05-tit-8 sub
\VUsaut{3.0mm}
\VUpk{10.2pt}{40}{வீட்டின் உள் அமைப்பு.}
\VUsaut{2.4mm}

%%K t05-apar-2 apar
{\centering\textit{(திட்டத்தைப் பார்க்க.)}\par}
\VUsaut{2.4mm}

%%K t05-01-48 p
\textsc{கட்டிடக் கலைஞர்}. --- வா இயோவான்னெஸ், வீட்டின் பல பகுதிகளையும்
உடனே உனக்குச் சுற்றிக் காட்டுகிறேன்.

%%K t05-01-49 p
\VUgras{தரைத் தளத்திற்குள்} \textsuperscript{(7)} நுழைவோம். இதோ
\VUgras{மணியின்} \textsuperscript{(11)} பொத்தான். மூன்று \VUgras{படிகளில்}
\textsuperscript{(14)} ஏறுகிறோம் ; அவை \VUgras{நுழைவுப் படிக்கட்டினுடையவை}
\textsuperscript{(9)}. பிறகு \VUgras{வாசற்படியைக்} \textsuperscript{(10)}
கடக்கிறோம் ; அது \VUgras{நுழைவுக் கதவினுடையது} \textsuperscript{(8)}.
இப்போது நாம் \VUgras{படிக்கட்டின்} \textsuperscript{(13)} அடியில்
இருக்கிறோம்.

%%K t05-01-50 p
\textsc{இயோவான்னெஸ்}. --- இது நடைபாதை.

%%K t05-01-51 p
\textsc{கட்டிடக் கலைஞர்}. --- இல்லை, இது \VUgras{நுழைவு அறை}
\textsuperscript{(12)}. \VUgras{நடைபாதை} \textsuperscript{(a)} நுனியில்
இருக்கிறது. வலப் பக்கத்தில் இதோ \VUgras{வரவேற்பு அறையும்}
\textsuperscript{(b)} \VUgras{உண்ணும் அறையும்} \textsuperscript{(c)} ;
உண்ணும் அறைக்கு எதிரே \VUgras{சமையலறையும்} \textsuperscript{(d)}
\VUgras{பணியாளர் அறையும்} \textsuperscript{(e)} ; பிறகு முன்னால்
\VUgras{எழுதும் அறை} \textsuperscript{(f)}. \VUgras{வராந்தா}
\textsuperscript{(23)} வரவேற்பு அறைக்கு அருகில் இருக்கிறது ; அதற்குத் தன்
\VUgras{கண்ணாடிக் கதவு} \textsuperscript{(25)} உண்டு.

%%K t05-01-52 p
இப்போது படிக்கட்டில் ஏறுவோம். \VUgras{கைப்பிடியைப்} \textsuperscript{(15)}
பிடித்துக் கொள், படிகளை எண்ணு. அவை பதினான்கு. இதோ \VUgras{படித் தளம்}
\textsuperscript{(m)} ; நாம் \VUgras{முதல் தளத்தில்} \textsuperscript{(27)}
இருக்கிறோம்.

%%K t05-tit-9 sub
\VUsaut{3.0mm}
\VUpk{10.2pt}{40}{முதல் தளம்.}
\VUsaut{3.0mm}

%%K t05-01-53 p
படிக்கட்டுக்கு எதிரே \VUgras{கழிப்பறையும்} \textsuperscript{(20)}
\VUgras{ஒப்பனை அறையும்} \textsuperscript{(j)} இருக்கின்றன. வலப் பக்கத்தில்
அழகான \VUgras{படுக்கை அறை} \textsuperscript{(g)} ; அதற்கு வீட்டின்
\VUgras{முகப்பில்} \textsuperscript{(1)} இரண்டு சாளரங்கள் உண்டு. இடப்
பக்கத்தில் \VUgras{குளியலறையும்} \textsuperscript{(1)}
\VUgras{விருந்தினர் அறையும்} \textsuperscript{(i)} ; அந்த அறைக்குப் பெரிய
\VUgras{படுக்கை மாடம்} \textsuperscript{(h)} உண்டு. படுக்கை அறைக்கு அருகில்
\VUgras{குழந்தைகள் அறை} \textsuperscript{(k)}. அது அகன்ற
\VUgras{மொட்டை மாடத்தை} \textsuperscript{(21)} ஒட்டி இருக்கிறது. இந்த
மாடத்திற்குக் கீழே இன்னொரு கழிப்பறை (20) இருக்கிறது ; சுவரில் இதோ
இரவு உணவுக்கு ஒலிக்கும் \VUgras{மணி} \textsuperscript{(22)}.

%%K t05-01-54 p
\textsc{இயோவான்னெஸ்}. --- \VUgras{இரண்டாம் தளத்திற்குப்} போவோம்.

%%K t05-01-55 p
\textsc{கட்டிடக் கலைஞர}. --- இரண்டாம் தளம் இல்லை. கூரைக்குக் கீழே
\VUgras{கூரை அறைகளும்} \textsuperscript{(33)} தானியக் கிடங்கும் (34)
மட்டுமே இருக்கின்றன. சுற்றிப் பார்த்து முடித்து விட்டோம், இறங்கலாம்.

%%K t05-01-56 p
\textsc{இயோவான்னெஸ்}. --- நன்றி ஐயா, இது மிகவும் சுவாரசியமாக இருந்தது.
நான் பார்த்த எல்லாவற்றையும் உடனே என் தந்தையிடம் சொல்லப் போகிறேன்.

\VUsaut{4.0mm}
\VUfilet{22mm}
